%!TEX root = ./../../thesis-trier.tex

In \Cref{sync_max:sec:dist_small_sets}
we give sufficient conditions for synchronizability in terms of different families of pairwise distinguishable
subsets of states.

\Cref{sync_max:sec:pairs_k_k-1} investigates families of subsets 
and their relations w.r.t.~distinguishability.

We show in \cref{sync_max:sec:shortestdist} that the set of distinguishing words of two subsets~$A,B$ 
of states of a given $n$-state automaton
is a regular language acceptable by an automaton of size $O(n^{|A|+|B|})$,
which also yields a bound on the length of shortest distinguishing words.

\Cref{sync_max:sec:completel-dist} introduces completely distinguishable automata.
We show that completely distinguishable automata with at least three states are synchronizing.
Then, we give several equivalent characterizations of this new class of automata.

Lastly, \cref{sync_max:sec:furtherexamples} gives examples of automata
where all pairs of distinct $2$-sets are distinguishable, but that are
not completely distinguishable.

\subsection{Distinguishability and Synchronizability}
\label{sync_max:sec:dist_small_sets}

%!TEX root = ../../thesis-trier.tex

In \cref{prelim:autsyncw:prop:two-sets-min-rank}
we prove, using \cref{sync_max:prop:two-sets-containing-2set},
that distinguishability of pairs of $k$-sets in automata with at least $k+1$ states 
implies synchronizability for $k \ge 3$, and for $k = 2$
existence of a word of rank at most two.
Then we list several consequences, for example,
that a strongly connected automaton in which pairs of $2$-sets
are distinguishable is synchronizing.
\Cref{sync_max:ex:2-sets-binary,sync_max:ex:2-sets}
give non-synchronizing automata where all pairs of $2$-sets
are distinguishable. Lastly, in~\cref{sync_max:prop:disjointklsets}, 
for automata whose transformation monoid contains a transitive group,
we give a sufficient condition for synchronizability in terms 
of the distinguishability of certain disjoint sets.

\begin{lemma}
	\label{sync_max:prop:two-sets-containing-2set}
	Let $\semiautomat{A}{Q}{\Sigma}{\delta}$.
	Suppose for all distinct states $p, q \in Q$
	there exist distinguishable subsets $A, B \subseteq Q$
	with $\{p,q\} \subseteq A \cap B$.
	Then $\mathcal A$ is synchronizing.	
\end{lemma}
\begin{proof}
	A word $u \in \Sigma^*$ that distinguishes $A$ from $B$
	maps at least one subset to a singleton subset,
	which implies $\applytrans{\delta}{p}{u} = \applytrans{\delta}{q}{u}$.
	As this holds for any pair of distinct states $p,q \in Q$, by \cref{prelim:thm:pair_sync_condition} 
	the automaton $\mathcal A$ is synchronizing.    			
\end{proof}

\begin{proposition} \label{prelim:autsyncw:prop:two-sets-min-rank} 	\label{sync_max:cor:k-ge-3-sync}
	Let $\semiautomat{A}{Q}{\Sigma}{\delta}$ and $2 \le k,l \le |Q|$ with $|Q| \ge k + 1$
	or $|Q| \ge l + 1$.
	Suppose each $k$-set is distinguishable from
	each distinct $l$-set of states in~$\mathcal A$
	and let $u$ be a word of minimal rank~$r$.
	For $k \ge 3$ or $l \ge 3$, we have $r = 1$, i.e., $\mathcal A$ is synchronizing.
	For $k = l = 2$, we have~$r \le 2$ where for $r = 2$
	the set $\delta(Q, u)$ is permuted by each letter, and so
	a union of sink components by \cref{prelim:autsyncw:lem:perm-each-letter-sink}.
\end{proposition}
\begin{proof}	
	Without loss of generality, let $k \le l$.
	First, suppose $l \ge 3$.
	Let $p,q \in Q$ be distinct.
	By assumption, there exist distinct $p_3, \ldots, p_k, p_{k+1} \in Q \setminus \{p,q\}$.
	Set $A = \{p,q,p_3, \ldots, p_k\}$.
	Next, there also exists an $l$-subset $B \subseteq Q$
	with $\{p,q,p_{k+1}\} \subseteq B$.
	Since $p_{k+1} \notin A$
	we have $A \ne B$.
	By assumption, $A$ and~$B$ are distinguishable.
	The, from \cref{sync_max:prop:two-sets-containing-2set} we conclude that $r = 1$. 
				
	Next, suppose $k = l = 2$.
	Assume $r \ge 3$. Then we can pick three distinct states $p,q,s \in \delta(Q, u)$.
	By assumption, the sets $\{p,q\}$ and $\{p,s\}$
	are distinguishable by some word $v \in \Sigma^*$.
	Since one of both sets is mapped to a singleton,
	we find $|\delta(Q, uv)| < |\delta(Q, u)|$, contradicting
	the choice of~$u$ as a word of minimal rank.
	So, we have $r \le 2$.
	Now, suppose $r = 2$.
	Assume there exists $a \in \Sigma$
	such that $\delta(\delta(Q, u), a) \ne \delta(Q, u)$.
	Since~$r$ is the minimal possible rank, $\delta(\delta(Q, u), a)$
	is a $2$-set.
	Then, by assumption, there exists a word $v$
	distinguishing $\delta(\delta(Q, u), a)$
	and $\delta(Q, u)$.
	However, then one of both subsets is mapped
	to a singleton, which is not possible
	since the minimal possible rank of a word is $2$.
	Hence, $\delta(\delta(Q, u), a) = \delta(Q, u)$
	for each $a \in \Sigma$.	
\end{proof}

Next, we list a few
consequences of \cref{sync_max:prop:two-sets-containing-2set,sync_max:cor:k-ge-3-sync}.

\begin{corollary} \label{sync_max:cor:two-sets-containing-2set}
	Let $\semiautomat{A}{Q}{\Sigma}{\delta}$ be an $n$-state automaton and $2 \le k < n$.
	
	\begin{enumerate}
		\item \label{sync_max:cor:two-sets-containing-2set:1} If $n \ge 4$ and all pairs of distinct $3$-sets
		are distinguishable in $\mathcal A$, then $\mathcal A$ is synchronizing.
		
		
		\item \label{sync_max:cor:two-sets-containing-2set:2} If for every $k$-set $A \subseteq Q$ there exists $B \subseteq Q$
		with $A \subseteq B$ such that $A$ and~$B$ are distinguishable,
		then $\mathcal A$ is synchronizing.
		
		\item \label{sync_max:cor:two-sets-containing-2set:item:kl}  Let $0 < l \le n - k$.
		If every $k$-set is distinguishable from 
		every $(k+l)$-set, then $\mathcal A$ is synchronizing.  
		
		\item \label{sync_max:cor:two-sets-containing-2set:4} If every $k$-set is distinguishable from the whole
		state set, then the automaton is synchronizing.
	\end{enumerate}	
\end{corollary}
\begin{proof}
	\ref{sync_max:cor:two-sets-containing-2set:1}: This follows from~\cref{sync_max:cor:k-ge-3-sync} with $k = 3$.
	
	\ref{sync_max:cor:two-sets-containing-2set:2}: For distinct $p,q \in Q$ choose some $k$-set with $p, q \in A$.
	By assumption, there exists $B \subseteq Q$ containing $A$ 
	so that $A$ and $B$ are distinguishable.
	Then \cref{sync_max:prop:two-sets-containing-2set} gives the claim.
	
	\ref{sync_max:cor:two-sets-containing-2set:item:kl} and~\ref{sync_max:cor:two-sets-containing-2set:4} clearly follow from~\ref{sync_max:cor:two-sets-containing-2set:2}.
\end{proof}

\begin{corollary} \label{prelim:aut-sync-words:subset:cor:no-letter-fixes-two-sets-and-dist-2-sets-implies-sync} \label{sync_max:prop:2-sets-dist-SCC-implies-sync}
Suppose $\semiautomat{A}{Q}{\Sigma}{\delta}$ is an automaton with at least three states in which all pairs of distinct $2$-sets
are distinguishable.
\begin{enumerate}
\item \label{sync_max:prop:2-sets-dist-SCC-implies-sync:1} If for distinct $p,q \in Q$ 
we have $\delta(\{p,q\}, a) \ne \{p,q\}$ for some $a \in \Sigma$,
then $\mathcal A$ is synchronizing.

\item \label{sync_max:prop:2-sets-dist-SCC-implies-sync:2} If~$\mathcal A$ is strongly connected, then it is synchronizing.

\item \label{sync_max:prop:2-sets-dist-SCC-implies-sync:3} Either $\mathcal A$ has a unique sink component and is synchronizing, or it has 
precisely two sink components each of size one.
\end{enumerate}
\end{corollary}
\begin{proof}
\ref{sync_max:prop:2-sets-dist-SCC-implies-sync:1}: Let $u \in \Sigma^*$ be a word of minimal possible rank.
By \cref{prelim:autsyncw:prop:two-sets-min-rank} 
we have $|\delta(Q, u)| \le 2$
where $|\delta(Q, u)| = 2$ is not possible
since otherwise every letter would permute this set
by \cref{prelim:autsyncw:prop:two-sets-min-rank}.

\ref{sync_max:prop:2-sets-dist-SCC-implies-sync:2}: Let~$u$ be a word of minimal rank
and assume $|\delta(Q, u)| \ge 2$.
By \cref{prelim:autsyncw:prop:two-sets-min-rank},
$\delta(Q, u)$ has size two and equals a union of sink components.
Since~$\mathcal A$ is strongly connected, $Q$ is its only sink component
and  we find $\delta(Q, u) = Q$,
which is not possible since~$|Q| \ge 3$.

\ref{sync_max:prop:2-sets-dist-SCC-implies-sync:3}:	By \cref{prelim:autsyncw:lem:minrank}, $\mathcal A$
has at most two sink components.
First, if it has a unique sink component~$S \subseteq Q$,
then $\mathcal A_{|S}$ is strongly connected
and all pairs of distinct $2$-sets from~$S$ are distinguishable.
Thus, by Item~\ref{sync_max:prop:2-sets-dist-SCC-implies-sync:2}, $\mathcal A_{|S}$ is synchronizing.
Since the set of maximal states equals the union of all sink components
by \cref{prelim:order:maximal-elements-upper-set}, $S$ equals the set of maximal states
and so~$\mathcal A$ is synchronizing by \cref{prelim:thm:SCC_sync_words}.
Next, suppose~$\mathcal A$ has two sink components.
Then it is not synchronizing by \cref{prelim:autsyncw:lem:minrank}.
By \cref{prelim:autsyncw:prop:two-sets-min-rank},
$\delta(Q, u)$ for a word of minimal rank~$u$
equals a union of sink components and has size two.
Since this set has non-empty intersection with both sink components by \cref{sync_max:eqn:upper_set}, 
both sink components must have size one.
\end{proof}

Next, \cref{sync_max:ex:2-sets-binary,sync_max:ex:2-sets} give infinite families of non-synchronizing automata 
where all pairs of $2$-sets are distinguishable, demonstrating that the case $k = 2$ in
\cref{prelim:autsyncw:prop:two-sets-min-rank} is special in the sense that
it does not imply synchronizability similar to the case $k > 2$.
By adding in these examples another letter that swaps the two maximal states and fixes every other state, we obtain non-synchronizing automata 
with a unique sink component where all pairs of $2$-sets are distinguishable.

\begin{example} 
	\label{sync_max:ex:2-sets}
	Set $\semiautomat{A_n}{[n]}{\Sigma}{\delta}$
	for $n \ge 3$
	where $\Sigma = \{ a_{ij} \mid i \in \{0,1\}, j \in \{2,\ldots,n-1\}\}$
	and, for $i \in \{0,1\}$, $j \in \{2,\ldots,n-1\}$,
	we have
	\[
	\delta(k, a_{ij}) = \begin{cases}
		i & \text{if } k = j, \\
		k & \text{otherwise,}
	\end{cases} 
	\]
	for $k \in [n]$.
	See \cref{sync_max:fig:2-sets}
	for~$\mathcal A_3$.
	Now, consider two distinct $2$-sets $\{p, q\}$
	and $\{s, t\}$ of states. First, without loss of generality,
	suppose $\{p, q\} \subseteq \{2,\ldots,n-1\}$.
	Then $\delta(p, u) = \delta(q, u) = 0$
	and $|\delta(\{s,t\}, u)| = 2$
	for $u = a_{0p} a_{0q}$.
	Otherwise, $\{0,1\}$ has non-empty intersection
	with both $2$-sets where we assume, without loss of 
	generality, $p = 0$ and $q \ne 1$.
	Then,
	$\delta(p, a_{0q}) = \delta(q, a_{0q}) = 0$
	and $|\delta(\{s, t\}, a_{0q})| = 2$.
\end{example}

\begin{example}
	\label{sync_max:ex:2-sets-binary}
	Let $\semiautomat{A_n}{Q_n}{\{a,b\}}{\delta_n}$ be the automaton given by the state diagram in \cref{sync_max:fig:2-sets-binary} where $n \ge 3$. Consider two distinct sets of states  $\{p_1, p_2\}, \{q_1, q_2\}$ with $p_1 < p_2$ and $q_1 < q_2$. %Note that this implies $p_1, q_1 \le n - 2$.
	First, without loss of generality, suppose $p_1 < q_1$.
	In this case, for $q_2 \le n - 2$, we have $\delta(\{q_1, q_2\}, a^{n - 2 - q_1}) = \{n - 2\}$ 
	and $|\delta(\{p_1, p_2\}, a^{n - 2 - q_1})| = 2$.
	For $q_2 = n - 1$ and $q_1 \ne n - 2$,
	we have $\delta(\{q_1, q_2\}, b^{n - 2 - q_1}) = \{n-1\}$
	and $|\delta(\{p_1, p_2\}, b^{n-2 - q_1})| = 2$.
	For $q_2 = n - 1$ and $q_1 = n - 2$,
	if $p_2 = n - 1$ the word $b^n$ distinguishes both subsets,
	and if $p_2 \le n - 2$ the word $a^n$ distinguishes both subsets.
	Next, suppose $p_1 = q_1$. In this case, without loss of generality,
	assume $p_2 < q_2$. % p_2 auf n - 2
	For $q_2 = n - 1$, we have $\delta(\{p_1, p_2\}, a^n) = \{n-2\}$
	and $\delta(\{q_1, q_2\}, a^n) = \{n - 1, n - 2\}$.
	Otherwise, for $q_2 < n - 1$,
	we have $\delta(\{p_1, p_2\}, u) = \{ n - 2 \}$
	and $\delta(\{q_1, q_2\}, u) = \{ n - 1, n - 2\}$
	for $u = a^{n - 3 - q_2}ba^n$. 
\end{example}

\todo[inline]{maßangaben auf em umstellen}
\begin{figure}
	\centering
	\begin{subfigure}[b]{0.44\textwidth}		
		\centering
		\begin{tikzpicture}[node distance = 2.5cm and 1.25cm]
			
			\node[state] (3) {$3$};
			
			\node[state, left = of 3] (4) {$4$};
			\node[state, right = of 3] (2) {$2$};
			\node[state, above = of 4] (0) {$0$};
			\node[state, above = of 2, yshift=-0.0cm] (1) {$1$};
			
			\path[->] (2) edge [bend left=10,above,sloped,pos=.7] node {$a_{02}$} (0)
			(2) edge [right] node {$a_{12}$} (1); 
			
			\path[->] (3) edge [bend left=15, above, sloped, pos=0.49] node {$a_{03}$} (0)
			(3) edge [bend right=15, right] node {$a_{13}$} (1)
			(4) edge [above, sloped] node {$a_{04}$} (0)
			(4) edge [bend right=15, pos=0.7,above, inner sep=9pt] node {$a_{14}$} (1);	
		\end{tikzpicture}	
		\caption[Non-Synchronizing and pairs of $2$-sets distinguishable]{Automaton as in \cref{sync_max:ex:2-sets}.}
		\label{sync_max:fig:2-sets}
	\end{subfigure}
	\hfill
	\begin{subfigure}[b]{0.55\textwidth}
		\centering
		\begin{tikzpicture}[inner sep=0pt, minimum size=20pt, node distance = 0.6cm and 0.7cm]
			\node[state] (0) {$0$};
			\node[state, right = of 0] (1) {$1$};
			\node[right = of 1] (dots) {$\cdots$};
			\node[state, right = of dots] (n-3) {$n-3$};
			\node[state, right = of n-3] (n-2) {$n-2$};
			\node[state, below = of n-2] (n-1) {$n-1$};
			
			\path[->] (0) edge [above] node {$a,b$} (1)
			(1) edge [above] node {$a,b$} (dots)
			(dots) edge [above] node {$a,b$} (n-3)
			(n-3) edge [above] node {$a$} (n-2)
			(n-3) edge [below] node {$b$} (n-1);
		\end{tikzpicture}
		\caption[Binary non-Synchronizing and pairs of $2$-sets distinguishable]{Binary automaton as in \cref{sync_max:ex:2-sets-binary}.} 
		\label{sync_max:fig:2-sets-binary}
	\end{subfigure}
	\caption[Tow non-synchronizing automata where all pairs of $2$-sets distinguishable]{Two non-synchronizing automata where all pairs of distinct $2$-sets are distinguishable, self-loops omitted.}
	\label{sync_max:fig:non-sync-two-sets-dist}
\end{figure}

Lastly, we consider the case that the transformation monoid contains a transitive group. In this case, we state a sufficient condition for synchronizability
in terms of the distinguishability of certain disjoint subsets. We need this result later in a more algebraic context.
For this reason, we give a formulation in terms of transformation monoids without reference to automata.

\begin{proposition} 
	\label{sync_max:prop:disjointklsets}
	Let $M \le_M \mathcal T_n$ be a transformation
	monoid with a transitive group of units
	and $2 \le k, l < n$ such that $kl < n$.
	If every $k$-set is distinguishable
	from every disjoint $l$-set,
	then $M$ contains a constant map. 		
\end{proposition}
\begin{proof}
	Without loss of generality, assume $k \le l$.
	Let $p, q \in [n]$ be two distinct states.
	Pick $A,B \subseteq [n]$ with $k = |A|$, $l = |B|$
	and $p,q \in A \cap B$ (note that the following argument includes the case $k = l = 2$ and $A = B = \{p,q\}$).
	As $G$ is transitive, by \cref{thm:separating_neumann} 
	we have $\applyfunc{A}{g} \cap B = \emptyset$
	for some $g \in G$.
	By assumption there exists $h \in M$
	mapping precisely one of both subsets $\applyfunc{A}{g}$
	or $B$ to a singleton subset.
	If $|\applyfunc{\applyfunc{A}{g}}{h}| = 1$,
	then, as $\{p,q\} \subseteq A$,
	the transformation $\concat{g}{h}$
	maps $\{p,q\}$ to a singleton.
	If $|\applyfunc{B}{h}| = 1$,
	then, as $\{p,q\} \subseteq B$,
	the transformation~$h$ maps $\{p,q\}$
	to a singleton.	
	Using the correspondence between automata
	and transformation semigroups (see \cref{prelim:remarks_trans_semigrp_semiautomata}),
	the pair criterion (\cref{prelim:thm:pair_sync_condition})
	implies that $M$ contains a constant map as all pairs of states $p,q$ 
	can be mapped to the same state.	
\end{proof}

\subsection{Classes of Pairwise Distinguishable Subsets}
\label{sync_max:sec:pairs_k_k-1} 

%!TEX root = ../../springer-mono-format/book.tex
%!TEX root = ../../thesis-trier.tex


\Cref{sync_max:prop:ksets} shows that distinguishability of $k$-sets for $k \ge 3$
in automata with at least $k+1$ states implies distinguishability of all $(k-1)$-sets.
Then \cref{aut-sync-words:dist-sets:prop:intersection-k-sets} 
investigates conditions equivalent to the distinguishability
of $(k+1)$-sets whose intersection has size~$k$.
\Cref{aut-syncw:subset:prop:kplusone} investigates
conditions when we can map, from two distinct subsets, one specific
set to a singleton, but not the other;
a property that implies distinguishability, but, in general, is stronger
since we can choose which subset is mapped to a singleton.
\Cref{sync_max:ex:2-3-sets-not-implying-2-sets-pre} 
gave an automaton having indistinguishable pairs
of disjoint $2$-sets, but where all $2$-sets
are distinguishable from sets of size at least three (this also implies distinguishability
of non-disjoint pairs of distinct $2$-sets by \cref{aut-syncw:subset:prop:kplusone}).
In \cref{syncmax:ex:disj-non-disj-two-sets} we give
two automata in which all disjoint $2$-sets are distinguishable,
but not all non-disjoint $2$-sets, and vice versa.
Thus, both properties are logically independent.
Our results use \cref{sync_max:lem:largersetsdist,sync_max:lem:sets-differ-one} stated first.

\begin{lemma} \label{sync_max:lem:largersetsdist}
	Let $\semiautomat{A}{Q}{\Sigma}{\delta}$
	be an automaton and $A, B, C \subseteq Q$ with $A \cap B \ne \emptyset$.
	Then, for a word~$u$
	with $\delta(A \cup C, u) = \{q\}$ for some $q \in Q$
	and $|\delta(B \cup C, u)| \ge 2$
	we have $|\delta(B, u)| \ge 2$.
\end{lemma}
\begin{proof}
	Since $A \cap B \ne \emptyset$
	we find $q \in \delta(B \cup C, u)$.
	We have $q \ne \delta(p, u)$ for some $p \in B \cup C$.
	Since $\delta(C, u) \subseteq \{q\}$
	we must have $p \in B$
	and so $\{q, \delta(p, u)\} \subseteq \delta(B, u)$.
\end{proof}

\begin{lemma} \label{sync_max:lem:sets-differ-one}
	Let $A,B \subseteq X$ and $p,q \in X$ such that $p \in A \setminus B$ and $q \in B \setminus A$.
	Then the following are equivalent.
 \begin{enumerate}
	\item \label{sync_max:lem:sets-differ-one:3a} $A \setminus \{p\} = B \setminus \{q\}$.
		
	\item \label{sync_max:lem:sets-differ-one:1a} $A \cup \{q\} = B \cup \{p\}$.
	
	\item \label{sync_max:lem:sets-differ-one:4a} $A \setminus B = \{p\}$ and $B \setminus A = \{q\}$.
	
	\item \label{sync_max:lem:sets-differ-one:5a} $A \cup B = A \cup \{q\}$ and $A \cap B = A \setminus \{p\}$.
		
	\item \label{sync_max:lem:sets-differ-one:2b} $A = (A \cap B) \cup \{p\}$
	and $B = (A \cap B)\cup \{q\}$.
	
	\item \label{sync_max:lem:sets-differ-one:2a} $A = ( B \setminus \{q\} ) \cup \{p\}$.	
\end{enumerate}
\end{lemma}
\begin{proof}
$\ref{sync_max:lem:sets-differ-one:3a}\Rightarrow\ref{sync_max:lem:sets-differ-one:1a}$: 
With $A \setminus \{p\} = B \setminus \{q\}$ we find
$
A \cup \{q\} = ( A \setminus \{p\} ) \cup \{p,q\}
= (B \setminus \{q\} ) \cup \{p,q\}
=  B \cup \{p\}.
$

$\ref{sync_max:lem:sets-differ-one:1a}\Rightarrow\ref{sync_max:lem:sets-differ-one:4a}$: 
Since $q \in B$ and $p \notin B$
we have $A \setminus B = (A \cup {q}) \setminus B = (B \cup \{p\}) \setminus B = \{p\}$.
Similarly, we find $B \setminus A = \{q\}$.

$\ref{sync_max:lem:sets-differ-one:4a}\Rightarrow\ref{sync_max:lem:sets-differ-one:5a}$: This follows since
$A \cup B = A \cup (B \setminus A) = A \cup \{q\}$
and $A \cap B = A \setminus (A \setminus B) = A \setminus \{p\}$.

$\ref{sync_max:lem:sets-differ-one:5a} \Rightarrow\ref{sync_max:lem:sets-differ-one:2b}$:
Since $A \cap B = A \setminus \{p\}$
we have $A = ( A \setminus \{p\} ) \cup \{p\} = (A \cap B) \cup \{p\}$.
Since $B \subseteq A \cup \{q\}$ and $q \in B \setminus A$,
we have $B \setminus A = \{q\}$.
This gives $B = (A \cap B) \cup (B \setminus A) = (A \cap B)\cup \{q\}$.

$\ref{sync_max:lem:sets-differ-one:2b} \Rightarrow\ref{sync_max:lem:sets-differ-one:2a}$:
Since $B = (A \cap B) \cup \{q\}$ and $q \notin A$ 
we have $B \setminus\{q\} = A \cap B$.
Thus $A = (A \cap B) \cup \{p\} = (B \setminus \{q\}) \cup \{p\}$.

$\ref{sync_max:lem:sets-differ-one:2a}\Rightarrow\ref{sync_max:lem:sets-differ-one:3a}$:
Since $p \notin B$, $A = ( B \setminus \{q\} ) \cup \{p\}$
implies $A \setminus \{p\} = B \setminus \{q\}$.
\end{proof}


\begin{proposition} \label{sync_max:prop:ksets}
	Let $\semiautomat{A}{Q}{\Sigma}{\delta}$
	be an automaton with at least $k+1$ 
	states for $k \ge 3$
	in which all pairs of distinct $k$-sets are distinguishable.
	Then pairs of distinct $(k-1)$-sets
	are distinguishable.
\end{proposition}
\begin{proof}
	Let $A,B \subseteq Q$ be two distinct
	$(k-1)$-sets. Then there exist
	$p \in A \setminus B$ and $q \in B \setminus A$.
	
	First, suppose $A \cup \{q\} = B \cup \{p\}$.
	In this case, by \cref{sync_max:lem:sets-differ-one},
	we find $A \cap B \ne \emptyset$
	and $|A \cup B| = k$.
	Since $|Q| \ge k + 1$
	there exists $s \in Q \setminus (A \cup B)$.
	Then $A \ne B$
	gives $A \cup \{s\} \ne B \cup \{s\}$
	and so, by assumption, there exists $u \in \Sigma^*$
	distinguishing these two $k$-sets.
	By \cref{sync_max:lem:largersetsdist}, $u$
	distinguishes also $A$ and~$B$.
	
	Next, suppose $A \cup \{q\} \ne B \cup \{p\}$.
	By assumption, there exists a word $u$ distinguishing
	$A  \cup \{q\}$ and $B  \cup \{p\}$.
	Without loss of generality, suppose $\delta(A \cup \{q\}, u) = \{s\}$
	for some $s \in Q$.
	Thus $\delta(q, u) = \delta(p, u) = s$
	and since $|\delta(B \cup \{p\}, u)| \ge 2$ 
	we have $\delta(r, u) \ne s$ for some $r \in B \setminus \{q\}$.
	Now, $q \in B$ and $r \in B$ 
	give $\{\delta(q, u), \delta(r, u) \} \subseteq \delta(B, u)$
	and so $|\delta(B, u)| \ge 2$.
\end{proof}

In automata with at most $k$ states, it is vacuously true
that all pairs of distinct $k$-sets are distinguishable.
Next, there exist automata (for example permutation automata)
with $k$ states for $k \ge 3$ where pairs of distinct $(k-1)$-sets
are indistinguishable.
Thus, we have to assume $|Q| \ge k + 1$ in \cref{sync_max:prop:ksets}.
Also, since two singleton sets are indistinguishable,
a similar statement as \cref{sync_max:prop:ksets} is false for $k = 2$.

\begin{proposition} \label{aut-sync-words:dist-sets:prop:intersection-k-sets}
	Let $\semiautomat{A}{Q}{\Sigma}{\delta}$ with 
	$|Q| \ge k+2$ for $k \ge 2$. Then the following are equivalent.
	
	\begin{enumerate}	
		\item \label{aut-sync-words:dist-sets:prop:intersection-k-sets:2} All pairs of $(k+1)$-sets
		whose intersection has size~$k$	are distinguishable.		
	
	\item \label{aut-sync-words:dist-sets:prop:intersection-k-sets:1} All pairs of $(l+1)$-sets
	whose intersection has size $1 \le l \le k$ are distinguishable.	
		
		\item \label{aut-sync-words:dist-sets:prop:intersection-k-sets:4} For distinct $A,B \subseteq Q$
		with $2 \le |A| \le k$ and $|B \setminus A| \ge 2$
		we have $|\delta(A, u)| = 1$ and $|\delta(B, u)| \ge 2$ for some $u \in \Sigma^*$,
		and pairs of distinct $2$-sets with non-empty intersection are distinguishable.
		
		\item \label{aut-sync-words:dist-sets:prop:intersection-k-sets:5} For a $k$-set $A$ and $2$-set $B$
		with $A\cap B = \emptyset$ there exists a word mapping $A$ to a singleton and~$B$ to a $2$-set,
		and all pairs of distinct $2$-sets with non-empty intersection are distinguishable.
		
		\item \label{aut-sync-words:dist-sets:prop:intersection-k-sets:3} For distinct $p_1, \ldots, p_{k+2} \in Q$
		there exists $u \in \Sigma^*$ such that $\delta(p_1, u) \ne \delta(p_{k+2}, u)$,
		and $\delta(p_1, u) = \cdots = \delta(p_{k+1})$	or $\delta(p_2, u) = \cdots = \delta(p_{k+2}, u)$.		
	\end{enumerate}
\end{proposition}
\begin{proof}
$\ref{aut-sync-words:dist-sets:prop:intersection-k-sets:2} \Rightarrow \ref{aut-sync-words:dist-sets:prop:intersection-k-sets:1}$: 
 Let $A,B \subseteq Q$ 
 be two $(l+1)$-sets with $|A \cap B| = l$
 for some $1 \le l \le k$.
 Then $A = (A \cap B) \cup \{p\}$ 
 and $B = (A \cap B) \cup \{q\}$
 for some $p \in A \setminus B$
 and some $q \in B \setminus A$.
 Thus we have $A \cup B = (A \cap B) \cup \{p,q\}$,
 which gives $|A\cup B| = l + 2$.
 Since $|Q| \ge k + 2$
 we can pick a $(k - l)$-set $C \subseteq Q \setminus ( A \cup B )$.
 Then $A \cup C$ and $B \cup C$ are two distinct $(k+1)$-sets
 whose intersection has size $k$.
 Since these sets are distinguishable by assumption, 
 by \cref{sync_max:lem:largersetsdist} we find that~$A$ and~$B$ are distinguishable.

$\ref{aut-sync-words:dist-sets:prop:intersection-k-sets:1} \Rightarrow \ref{aut-sync-words:dist-sets:prop:intersection-k-sets:4}$: For $l = 1$
 we find distinguishability of
 distinct $2$-sets with non-empty intersection.
 Next, let~$A,B \subseteq Q$ be as in the statement and $p,q \in B \setminus A$ be distinct.
 The sets $A \cup \{p\}$
 and $A \cup \{q\}$ are distinct with intersection~$A$
 of size at most~$k$.
 So, by assumption, distinguishable by some word~$u \in \Sigma^*$.
 This yields $|\delta(A, u)| = 1$
 and $\delta(p, u) \ne \delta(q, u)$
 and so $|\delta(B, u)| \ge 2$.

$\ref{aut-sync-words:dist-sets:prop:intersection-k-sets:4} \Rightarrow \ref{aut-sync-words:dist-sets:prop:intersection-k-sets:5}$: This is clear.

$\ref{aut-sync-words:dist-sets:prop:intersection-k-sets:5} \Rightarrow \ref{aut-sync-words:dist-sets:prop:intersection-k-sets:3}$: Let $p_1, \ldots, p_{k+2} \in Q$ be distinct.
 By considering $A = \{p_2, \ldots, p_{k+1}\}$ and $B = \{p_1, p_{k+2}\}$
 we find that there exists $u \in \Sigma^*$
 such that $\delta(p_2, u) = \cdots = \delta(p_{k+1}, u)$
 and $\delta(p_1, u) \ne \delta(p_{k+2}, u)$.
 For $\delta(p_2, u) \in \delta(B, u)$
 we obtain the claim.
 Otherwise, for $\delta(p_2, u) \notin \delta(B, u)$
 we have distinct $2$-sets $\{\delta(p_1, u), \delta(p_2, u)\}$
 and $\{\delta(p_{k+2}, u), \delta(p_2, u)\}$, which
 are distinguishable by some word~$v \in \Sigma^*$.
 In this case the word $uv$ fulfills the claim.	
 
$\ref{aut-sync-words:dist-sets:prop:intersection-k-sets:3} \Rightarrow \ref{aut-sync-words:dist-sets:prop:intersection-k-sets:2}$: Given two distinct $(k+1)$-sets $A,B \subseteq Q$
with $|A \cap B| = k$,
we can apply the assumption to $A \cap B = \{p_2, \ldots, p_{k+1}\}$, $p_1 \in A \setminus (A \cap B)$,
and $p_{k+2} \in B \setminus (A \cap B)$
and obtain that $A$ and~$B$ are distinguishable.
\end{proof}

\begin{proposition} \label{aut-syncw:subset:prop:kplusone}
	Let $\semiautomat{A}{Q}{\Sigma}{\delta}$ with $|Q| > k \ge 2$. Then the following are equivalent.	
	\begin{enumerate}		
		\item  \label{aut-syncw:subset:prop:kplusone:1}  Two distinct $A,B \subseteq Q$ with $2 \le |A| \le k$ and $A \cap B \ne \emptyset$ are distinguishable.		
		
		\item  \label{aut-syncw:subset:prop:kplusone:2} Two distinct $A,B \subseteq Q$ with $|A| = k$ and $A \cap B \ne \emptyset$ are distinguishable.
		
		\item  \label{aut-syncw:subset:prop:kplusone:3} For a $k$-set $A \subseteq Q$ and $p \notin Q \setminus A$ the sets $A$ and $A \cup \{p\}$ are distinguishable.
		
		\item  \label{aut-syncw:subset:prop:kplusone:4} For distinct $k$-sets $A,B \subseteq Q$ with $A \cap B \ne \emptyset$ we have $|\delta(A, u)| = 1$ and $|\delta(B, u)| \ge 2$ for some $u \in \Sigma^*$.
		
		\item  \label{aut-syncw:subset:prop:kplusone:5} For $A \subseteq Q$ with $2 \le |A| \le k$ and $p \notin Q \setminus A$ the sets $A$ and $A \cup \{p\}$ are distinguishable.
				
		\item  \label{aut-syncw:subset:prop:kplusone:6} For $A,B \subseteq Q$ with $|A|,|B| \ge 2$ and $B\not\subseteq A$
		where $|A| \le k$ and $|A| = k$ implies $A \cap B \ne \emptyset$
		we have $|\delta(A, u)| = 1$ and $|\delta(B, u)| \ge 2$ for some $u \in \Sigma^*$.		
	\end{enumerate}
\end{proposition}
\begin{proof}
	$\ref{aut-syncw:subset:prop:kplusone:1} \Rightarrow \ref{aut-syncw:subset:prop:kplusone:2}$: This is clear.	
	
	$\ref{aut-syncw:subset:prop:kplusone:2} \Rightarrow \ref{aut-syncw:subset:prop:kplusone:3}$: This follows from the assumption with $A \cup \{p\}$ in the role of~$B$.
	
	$\ref{aut-syncw:subset:prop:kplusone:3} \Rightarrow \ref{aut-syncw:subset:prop:kplusone:4}$: Since $|A| = |B|$
	there exists $p \in B \setminus A$.
	Then a word distinguishing $A$ and $A \cup \{p\}$
	maps $A$ to a singleton and $B$ to a set of size at least two
	since $A \cap B \ne \emptyset$ and $p \in B$.
	
	$\ref{aut-syncw:subset:prop:kplusone:4} \Rightarrow \ref{aut-syncw:subset:prop:kplusone:5}$: Since $|Q| > k$
	there exists $C \subseteq Q \setminus \{p\}$ such that $A \cup C$ is a $k$-set.
	Set $D = ( ( A \cup C ) \setminus \{q\} ) \cup \{p\}$ for some $q \in A \cup C$.	
	By assumption, there exists $u$ such that $\delta(A \cup C, u) = \{s\}$
	and $|\delta(D, u)| \ge 2$. Then we must have $\delta(D, u) = \{s, \delta(p, u)\}$.
	Since $A \cap D \ne \emptyset$,
	this gives $\delta(A \cup \{p\}, u) = \{s, \delta(p, u)\}$.
		
	$\ref{aut-syncw:subset:prop:kplusone:5} \Rightarrow \ref{aut-syncw:subset:prop:kplusone:6}$: By assumption, 
	there exist $p \in B \setminus A$ and $q \in B$ with $p \ne q$
	such that $|A \cup\{q\}| \le k$.	
	We find $u$ distinguishing $A \cup \{q\}$ and $A \cup \{q,p\}$.
	This word must map~$A \cup \{q\}$ to a singleton and $p$ to a different state than those in~$A \cup \{q\}$.
	Since $p,q \in B$, it maps $B$ to a set of size at least two.	
	
	$\ref{aut-syncw:subset:prop:kplusone:6} \Rightarrow \ref{aut-syncw:subset:prop:kplusone:1}$: By assumption, 
	there exists $u$ mapping $A$ to a singleton and $A \cup B$ to a set of size at least two.
	Since $A \cap B \ne \emptyset$, this word also maps~$B$ to a set of size at least two.
\end{proof}

Note that \crefitem{aut-syncw:subset:prop:kplusone}{aut-syncw:subset:prop:kplusone:3}
implies that for non-disjoint $2$-sets $A,B$
there exist two distinguishing words, one mapping $A$ to a singleton
and another mapping~$B$ to a singleton, and for $k \ge 3$
by \crefitem{aut-syncw:subset:prop:kplusone}{aut-syncw:subset:prop:kplusone:2},
that for every pair of distinct $2$-sets $A,B$
there exist two distinguishing words, each mapping a different
set to a singleton.

For the automaton from \cref{sync_max:ex:2-3-sets-not-implying-2-sets-pre},
the conditions from \cref{aut-syncw:subset:prop:kplusone} hold
for $k = 2$.

\begin{example} \label{syncmax:ex:disj-non-disj-two-sets} \Cref{sync_max:fig:2-sets-disjoint-vs-nondisjoint} 
	gives two examples showing that neither distinguishability of disjoint pairs of $2$-subsets
	nor distinguishable of distinct non-disjoints pairs of $2$-subsets 
	is sufficient for distinguishability of all pairs of $2$-subsets.
	\Cref{sync_max:fig:2-3-sets-not-implying-2-sets} gives
	another example where non-disjoint pairs of $2$-sets are distinguishable
	but disjoint pairs of $2$-sets are indistinguishable.
\end{example} 

\begin{figure}[ht]
	\centering
	\begin{subfigure}[t]{0.5\textwidth}
		\centering
		\begin{tikzpicture}[node distance = 1cm and 1cm]
			\node[state] (0) {$0$};
			\node[state, below left  = of 0] (1) {$1$};
			\node[state, below right = of 0] (2) {$2$};
			\node[state, right       = of 0] (3) {$3$};
			
			\path[->] (1) edge [left] node {$a$} (0)
			(2) edge [right] node {$a$} (0)
			(0) edge [loop left] node {$a$} (0)
			(3) edge [loop right] node {$a$} (3);
		\end{tikzpicture}	
		\caption{All disjoint pairs of $2$-sets are distinguishable, but $\{0,1\}$ is not distinguishable from $\{0,2\}$.}
	\end{subfigure}%
	\hfill
	\begin{subfigure}[t]{0.45\textwidth}
		\centering
		\begin{tikzpicture}[node distance = 1.1cm and 1.5cm]
			\node[state] (0) {$0$};
			\node[state, right = of 0] (1) {$1$};
			\node[state, below = of 1] (2) {$2$};
			\node[state, left  = of 2] (3) {$3$};
			
			\path[->] (1) edge [bend left,above] node {$a$} (0)
			(0) edge [loop left] node {$a$} (0)
			(2) edge [above] node {$a,b$} (3)
			(3) edge [loop left] node {$a$} (3);
			
			\path[->] (0) edge [bend left,above] node {$b$} (1)
			(1) edge [right] node {$b$} (2)
			(3) edge [left]  node {$b$} (0);
		\end{tikzpicture}	
		\caption{Pairs of distinct non-disjoint $2$-sets are distinguishable and pairs of
			disjoint~$2$-sets are indistinguishable.}
	\end{subfigure}		
	\caption[Relations between disjoint and non-disjoint pairs of subsets w.r.t.~distinguishability]{Two 
		automata showing that the properties 
		of distinguishability of disjoint and non-disjoint $2$-sets do not imply each other, see \cref{syncmax:ex:disj-non-disj-two-sets}.		
		
	}	
	\label{sync_max:fig:2-sets-disjoint-vs-nondisjoint}
\end{figure}
%
% Noch Fall disjointer k-sets. Was dazu äquivalent?

\subsection{The Set of Words Distinguishing Two Sets and a Length Bound} %Shortest Distinguishing Words} 
\label{sync_max:sec:shortestdist} 

%!TEX root = ../../thesis-trier.tex

Let $A$ and $B$ be two subsets of states in an $n$-state automaton.
Here, we show that the set of words distinguishing $A$ from $B$
is acceptable by an automaton with $O(n^{|A|+|B|})$ states.
As a consequence, a shortest distinguishing word has length in~$O(n^{|A|+|B|})$.

\begin{theorem}
	\label{sync_max:thm:distinguishing_word_bound}
	Let $\semiautomat{A}{Q}{\Sigma}{\delta}$ 
	be an $n$-state automaton and $A, B \subseteq Q$
	with $|A|, |B| \ge 2$.	
	Then the set of words distinguishing $A$ from $B$ is regular
	and acceptable by an automaton with 
	\[
	\left( n + \binom{n}{2} + \binom{n}{3} + \cdots + \binom{n}{|A|} \right) \cdot
	\left( n + \binom{n}{2} + \binom{n}{3} + \cdots + \binom{n}{|B|} \right) - n + 1
	\]
	states.	Furthermore, for a shortest distinguishing word $u \in \Sigma^*$ we have
	\[
	|u| \le \left( \binom{n}{2} + \binom{n}{3} + \cdots + \binom{n}{|A|} \right) \cdot
	\left( \binom{n}{2} + \binom{n}{3} + \cdots + \binom{n}{|B|} \right)
	\]
\end{theorem}
\begin{proof}
	Let  $\automat{B}{Q'}{\Sigma}{\mu}{(A,B)}{E}$ where
	\begin{align*}
		Q' & = \{ R \subseteq Q : 0 < |R| \le |A| \} \times \{ P \subseteq Q : 0 < |P| \le |B| \}, \\
		E & = \{ (R, P) \in Q' \times Q' : |R| = 1 \text{ if and only if } |P| > 1 \}, \\
		\applytrans{\mu}{(R, P)}{x} & = (\applytrans{\delta}{R}{x}, \applytrans{\delta}{P}{x}) \text{ for } (R, P) \in Q' \times Q' \text{ and } x \in \Sigma.
	\end{align*}
    Set $F =  \{ \{q\} \mid q \in Q \}$.
    Then $\mathcal B / {\sim_{F \times F}}$ with start state $\pi_{\sim_{F \times F}}((A,B))$
    and final state set $\pi_{\sim_{F\times F}}(E)$
    accepts the set of words distinguishing $A$ from $B$.
    This automaton has the claimed number of states.
        
    Next, let $u \in \Sigma^*$ be a shortest distinguishing word.
    We have $|u| > 0$ and thus we can write $u = wx$ for some $w \in \Sigma^*$ and $x \in \Sigma$.
    By the minimal choice and since~$u$ is a distinguishing word, 
    for each proper prefix $v \in \Sigma^*$
    of~$u$, for $(R, P) = \mu(A, B, v)$
    we have $(R, P) \notin E$
    with $|R| \ge 2$ and $|P| \ge 2$.    
    Thus, since we have 
    \[
    |\{ (R, P) \in Q' \times Q' : |R| \ge 2, |P| \ge 2 \}| = \left( \binom{n}{2} + \binom{n}{3} + \cdots + \binom{n}{|A|} \right) \cdot 
    \left( \binom{n}{2} + \binom{n}{3} + \cdots + \binom{n}{|B|} \right)
    \]
    and all states traverses by prefixes of~$w$ are distinct,
    we find
    \[
    |w| \le \left( \binom{n}{2} + \binom{n}{3} + \cdots + \binom{n}{|A|} \right) \cdot 
    \left( \binom{n}{2} + \binom{n}{3} + \cdots + \binom{n}{|B|} \right) - 1
    \]
    This gives the claimed bound for~$u$.
\end{proof}

\subsection{Completely (Subset-Pair) Distinguishable Automata}
\label{sync_max:sec:completel-dist} 

%!TEX root = ./../../springer-mono-format/book.tex
%!TEX root = ./../../thesis-trier.tex

\IN{completely distinguishable!automaton}
An automaton $\semiautomat{A}{Q}{\Sigma}{\delta}$ is called \emph{completely subset-pair distinguishable} (or just \emph{completely distinguishable} for short) 
if every two distinct non-empty and non-singleton subsets $A,B \subseteq Q$ are distinguishable in~$\mathcal A$.

Recall that $L(\mathcal P(\mathcal A)_{Q, F}) = \Syn(\mathcal A)$
with $F = \{\{q\} \mid q \in Q\}$, see~\cref{sync_max:sec:power_automat_for_SynA}. 
By \crefitem{prelim:thm:char_minaut}{thm:minaut:item:statesdist},
distinguishability of pairs of distinct states is a necessary condition 
for minimality.
So, applying this to $\mathcal P(\mathcal A)_{Q, F}$,
complete distinguishability
is a necessary condition for the minimal automaton
of $\Syn(\mathcal A)$ to have $2^{|Q|} - |Q|$ states.

First, in \cref{sync_max:thm:cerny_compl_dist}, as an example,
we prove that the \Cerny automata from \cref{intro:fig:cerny}
are completely distinguishable (this also follows from \cref{sync_max:thm:circular_sync_max} given later,
but as an example we give a self-contained proof here).
Then, in \cref{sync_max:prop:dist_sync},
we show that completely distinguishable automata
with at least three states are synchronizing.
\Cref{sync_max:thm:char_distinguishable} collects various equivalent characterizations,
and \cref{sync_max:cor:sufficient-compl-dist,cor:compl-dist-transitive}
are two consequences of these characterizations.
\Cref{sync_max:ex:compl-dist-non-cr,sync_max:ex:compl-dist-non-cr2,syncmax:ex:compl-dist-not-SCC} 
give completely distinguishable automata that are not completely reachable.
\Cref{syncmax:prop:binary-rank-two-letter} 
proves that in binary completely distinguishable automata where one letter has rank two,
the other letter cyclically permutes the states.
\cref{sync_max:prop:dist-n-1-set} lists an equivalence for a set missing one state being
distinguishable from the whole state set (a condition that appears in \cref{sync_max:thm:char_distinguishable}
for each set missing one state).

\begin{theorem}
	\label{sync_max:thm:cerny_compl_dist}
	The \Cerny automata $\mathcal C_n = ([n], \{a,b\}, \delta_n)$ from \cref{intro:fig:cerny}
	are completely distinguishable.
\end{theorem}
\begin{proof}
For $i \in [n]$ 
set $u_i = b^{n - i} a b^i$.
Then, for $j \in [n]$,
\begin{equation} \label{sync_max:thm:cerny_compl_dist:eqn}
\applytrans{\delta_n}{j}{u_i} = \begin{cases}
	(j + 1 ) \bmod n & \text{if } i = j, \\
	j             & \text{if } i \ne j.	
\end{cases}
\end{equation}
Now, let $A, B \subseteq [n]$ be distinct with $|A|, |B| \ge 2$. %two distinct non-empty and non-singleton subsets. 
Without loss of generality, 
suppose $i \in A \setminus B$. Set $u = u_{i+1} \cdots u_0 u_1 \cdots u_{i-1}$.
Then $\delta(i, u) = i$, $\applytrans{\delta}{B}{u} \subseteq \applytrans{\delta}{[n] \setminus \{i\}}{u} = \{ i - 1 \}$
and $\applytrans{\delta}{A}{u} = \{i,i-1\}$. So,~$u$
distinguishes~$A$ from~$B$.
\end{proof} % then u = ... disnigushes a from B as ...

\crefitem{sync_max:cor:two-sets-containing-2set}{sync_max:cor:two-sets-containing-2set:item:kl} yields
that completely distinguishable automata with at least three states are synchronizing.

\begin{proposition}
	\label{sync_max:prop:dist_sync}
	Let $\semiautomat{A}{Q}{\Sigma}{\delta}$ 
	be a completely distinguishable automaton
	with $|Q| \ge 3$. Then $\mathcal A$
	is synchronizing.
\end{proposition} 
\begin{proof}
	This follows from \crefitem{sync_max:cor:two-sets-containing-2set}{sync_max:cor:two-sets-containing-2set:item:kl},
	since every $2$-set is distinguishable from every $3$-set.
\end{proof}

Since the empty set is distinguishable from every other set if and only if 
the automaton is synchronizing, by \cref{sync_max:prop:dist_sync}
an automaton with at least three states is completely distinguishable if every pair 
of non-singleton subsets is distinguishable.

\begin{remark}\label{syncmax:rem:two-state-compl-dist}
	An automaton with at most two states admits no pair of distinct non-empty
	and non-singleton subsets of states. Thus, it is
	completely distinguishable, even if it is non-synchronizing.	
\end{remark}

Next, we formulate various equivalent characterizations.

\begin{theorem}
	\label{sync_max:thm:char_distinguishable}
	Let $\semiautomat{A}{Q}{\Sigma}{\delta}$ be an $n$-state automaton
	with $n \ge 3$.
	The following are equivalent:
	\begin{enumerate}
		\item \label{dist:item:all_pairs}
		$\mathcal A$ is completely distinguishable.
		
		\item \label{dist:item:subsets} 
		For every $A, B \subseteq Q$
		with $A \subsetneq B$ and $A$ non-empty
		there exists $u \in \Sigma^*$
		such that $1 = |\applytrans{\delta}{A}{u}| < |\applytrans{\delta}{B}{u}|$.
		
		\item \label{dist:item:subset_single_point} For every non-empty $A \subseteq Q$ and $q \in Q \setminus A$ there exists $u \in \Sigma^*$ 
		such that $\applytrans{\delta}{q}{u} \notin \applytrans{\delta}{A}{u}$
		and $|\applytrans{\delta}{A}{u}| = 1$.		
						
		\item \label{dist:item:single_point} For every $q \in Q$ there exists $u \in \Sigma^*$ such that
		$1 = |\applytrans{\delta}{Q \setminus \{q\}}{u}|$ and $\applytrans{\delta}{q}{u} \notin \applytrans{\delta}{Q \setminus \{q\}}{u}$, i.e., $Q$ is distinguishable from every $(n-1)$-set.
		
		\item \label{dist:item:one-two-set} For two distinct non-empty and non-singleton subsets
		there exists a word mapping one to a singleton and the other to a $2$-set.
					
		\item \label{dist:item:pqr-r-one-lane} For distinct $p,q,r \in Q$ there exists $u \in \Sigma^*$ 
		such that $\delta(p, u) = \delta(q, u)$ and 
		\[
		\delta^{-1}(\{\delta(r,u)\}, u) = \{r\},
		\] i.e.,
		$\delta(s, u) \ne \delta(r, u)$ for all $s \in Q \setminus \{r\}$.
		
		\item \label{dist:item:compress-two-sets} For every $A, B \subseteq Q$ with $\min\{|A|, |B|\} \ge 2$
		and $A \subsetneq B$ there exists $u \in \Sigma^*$ such that $|\delta(A, u)| < |\delta(B, u)| < |B|$.		
				
		\item \label{dist:item:compress} For every $A \subseteq Q$
		with $|A| \ge 2$ and $p \in Q \setminus A$
		there exists $u \in \Sigma^*$
		such that $\applytrans{\delta}{p}{u} \notin \applytrans{\delta}{A}{u}$
		and $|\applytrans{\delta}{A}{u}| < |A|$. 
				
		\item \label{dist:item:dcps-nondisjoint} For every distinct non-empty and non-singleton $A, B \subseteq Q$ with $A \cap B \ne \emptyset$
		there exists $u \in \Sigma^*$
		with $3 \le |\applytrans{\delta}{A}{u}| + |\applytrans{\delta}{B}{u}| < |A| + |B|$
		and $\applytrans{\delta}{A}{u} \ne \applytrans{\delta}{B}{u}$.		
		
		\item \label{dist:item:dcps} For every distinct non-empty and non-singleton $A, B \subseteq Q$  
		there exists $u \in \Sigma^*$
		with $3 \le |\applytrans{\delta}{A}{u}| + |\applytrans{\delta}{B}{u}| < |A| + |B|$
		and $\applytrans{\delta}{A}{u} \ne \applytrans{\delta}{B}{u}$.		
	\end{enumerate}
\end{theorem}
\begin{proof}
	$\ref{dist:item:all_pairs} \Rightarrow \ref{dist:item:subsets}$:
	By assumption there exists $u \in \Sigma^*$
	distinguishing $A$ from $B$.
	As $A \subsetneq B$
	we must have $|\applytrans{\delta}{A}{u}| = 1$
	and $|\applytrans{\delta}{B}{u}| > 1$.
	
	$\ref{dist:item:subsets} \Rightarrow \ref{dist:item:subset_single_point}$:
	For $A \subseteq Q$ and $q \in Q \setminus A$ the claim follows by applying Item~\ref{dist:item:subsets} to~$A$ and $B = A \cup \{q\}$.	
	
	$\ref{dist:item:subset_single_point} \Rightarrow \ref{dist:item:single_point}$: 
	For $q \in Q$ the claim follows	by applying Item~\ref{dist:item:subset_single_point} to $A = Q \setminus \{q\}$ and $q \in Q \setminus A$.
		
	$\ref{dist:item:single_point} \Rightarrow \ref{dist:item:one-two-set}$:
	Let $A,B \subseteq Q$ be two distinct subsets of size at least two.
	Without loss of generality, suppose $q \in A \setminus B$.
	By assumption, there exists $u \in \Sigma^*$
	such that $|\delta(Q \setminus \{q\}, u)| = 1$
	and $\delta(q, u) \notin \delta(Q \setminus \{q\})$.
	Since $B \subseteq Q \setminus\{q\}$
	and $A$ contains some state distinct from~$q$,    
	we have $|\delta(B, u)| = 1$
	and $|\delta(A, u)| = 2$.	
	
	$\ref{dist:item:one-two-set} \Rightarrow \ref{dist:item:pqr-r-one-lane}$:
	Let $p,q,r\in Q$ be distinct. Consider
	$Q \setminus \{r\}$ and $Q$.
	By assumption, there exists a word $u \in \Sigma^*$ mapping one to a singleton
	and the other to a $2$-set,
	where we must have $|\delta(Q \setminus \{r\}, u)| = 1$
	and $|\delta(Q, u)| = 2$, which gives $\delta(r, u) \notin \delta(Q \setminus \{r\}, u)$.
	Then as $p,q \in Q \setminus\{r\}$ we have $\delta(p, u) = \delta(q, u)$
	and, as for $s \in Q \setminus \{r\}$ we have $\delta(s, u) = \delta(p, u)$,
	we have $\delta^{-1}(\{\delta(r,u)\}, u) = \{ s \in Q \mid \delta(s, u) = \delta(r, u)\} = \{ r \}$, i.e.,
	$r$ is the only state mapped to $\delta(r, u)$.	
		
	$\ref{dist:item:pqr-r-one-lane} \Rightarrow \ref{dist:item:compress-two-sets}$: Pick distinct states $p,q \in A$ and $r \in B \setminus A$.
	By assumption there exists $u \in \Sigma^*$
	such that $\delta(p, u) = \delta(q, u)$
	and $\delta(r, u) \ne \delta(s, u)$ for all $s \in Q \setminus \{r\}$.
	So $\delta(r, u) \notin \delta(A, u)$
	and thus $|\delta(A, u)| < |\delta(B, u)|$.
	Furthermore, since $\{p,q\} \subseteq B$
	we have $|\delta(B, u)| < |B|$.
	
	$\ref{dist:item:compress-two-sets} \Rightarrow \ref{dist:item:compress}$: 
	Consider~$A$ and $B = A \cup \{p\}$.
	By assumption, we have $|\delta(A, u)| < |\delta(B, u)| < |B|$
	for some $u \in \Sigma^*$.
	Since $|\delta(A, u)| < |\delta(B, u)|$,
	we find $\delta(p, u) \notin \delta(A, u)$.	
	We have $|B| = |A| + 1$ and $\delta(B, u) = \delta(A, u) \cup \delta(\{p\}, u)$
	so that $|\delta(B, u)| = |\delta(A, u)| + 1$.
	Thus,  $|\delta(B, u)| < |B|$
	yields $|\delta(A, u)| < |A|$.
			
	$\ref{dist:item:compress} \Rightarrow \ref{dist:item:dcps-nondisjoint}$:
	Let $A, B \subseteq Q$ be distinct with $|A|, |B| \ge 2$ and $A \cap B \ne \emptyset$.
	Without loss of generality, suppose $p \in B \setminus A$ for some $p \in Q$.
	By assumption we have $|\delta(A, u)| < |A|$
	and $\delta(p, u) \notin \delta(A, u)$ for some $u \in \Sigma^*$.
	Then we get $|\applytrans{\delta}{A}{u}| + |\applytrans{\delta}{B}{u}| < |A| + |B|$
	and, as $\delta(p, u) \in \delta(B, u)$, we get
	$\delta(A, u) \ne \delta(B, u)$.
	Furthermore, $\delta(A \cap B, u) \cup \{\delta(p,u)\} \subseteq \delta(B, u)$ with $A \cap B \ne \emptyset$ and $\delta(p, u) \notin \delta(A \cap B, u)$ gives $|\delta(B, u)| \ge 2$ so that $|\applytrans{\delta}{A}{u}| + |\applytrans{\delta}{B}{u}| \ge 3$.
	
	$\ref{dist:item:dcps-nondisjoint} \Rightarrow \ref{dist:item:dcps}$: Let $A, B \subseteq Q$ be distinct with $|A|, |B| \ge 2$. Without loss of generality, let $p \in B \setminus A$ and $q \in B \setminus\{p\}$. By the assumption applied to the distinct non-empty and non-singleton sets $A \cup \{q\}$ and $A \cup \{p,q\}$ there exists $u \in \Sigma^*$ with $3 \le |\delta(A \cup \{q\}, u)| + |\delta(A \cup \{p, q\}, u)| < 2|A|+3$ (where we even have 
	$|\delta(A \cup \{q\}, u)| + |\delta(A \cup \{p, q\}, u)| < 2|A| + 1$ for $q \in A$) and $\delta(A \cup \{q\}, u) \ne \delta(A \cup \{p, q\}, u)$. Since $\delta(A \cup \{p, q\}, u) = \delta(A \cup \{q\}, u) \cup \delta(\{p\}, u)$ we find $\delta(p, u) \notin \delta(A \cup \{q\}, u)$.
	Hence, as $\delta(p, u) \in \delta(B, u)$, we have $\delta(A, u) \ne \delta(B, u)$ 
	and $|\delta(B, u)| \ge 2$ as $\{\delta(p,u), \delta(q, u)\} \subseteq \delta(B, u)$
	where $\delta(p, u) \ne \delta(q, u)$,	
	which gives $3 \le |\delta(A, u)| + |\delta(B, u)|$. Next, for $\delta(\{p,q\}, u) \cap \delta(A, u) = \emptyset$ we have 
	\begin{align*}
		2|\delta(A, u)| + 3 & = |\delta(A, u)| + |\delta(\{q\}, u)| + |\delta(A, u)| + |\delta(\{p,q\}, u)| \\
		& = |\delta(A \cup \{q\}, u)| + |\delta(A \cup \{p, q\}, u)| \\
		& < 2|A| + 3.
	\end{align*}
	Now, $2|\delta(A, u)| + 3 < 2|A| + 3$ gives $|\delta(A, u)| < |A|$. So, we find $|\delta(A, u)| + |\delta(B, u)| < |A| + |B|$. Otherwise, suppose $\delta(\{p,q\}, u) \cap \delta(A, u) \ne \emptyset$, which gives $\delta(A, u) \cap \delta(B, u) \ne \emptyset$. In this case, either we have $|\delta(A, u)| + |\delta(B, u)| < |A| + |B|$ and the claim holds with~$u$, or $|\delta(A, u)| = |A| \ge 2$ and $|\delta(B, u)| = |B| \ge 2$, and by the assumption applied to $\delta(A, u)$ and $\delta(B, u)$ there exists $v \in \Sigma^*$ with $3 \le |\delta(A, uv)| + |\delta(B, uv)| < |\delta(A, u)| + |\delta(B, u)|$ and $\delta(A, uv) \ne \delta(B, uv)$. Then the claim holds with~$uv$ as $|\delta(A, u)| + |\delta(B, u)| \le |A| + |B|$.
	
	$\ref{dist:item:dcps} \Rightarrow \ref{dist:item:all_pairs}$:
	Let $A,B \subseteq Q$
	be distinct with $|A|, |B| \ge 2$.
	We do induction on $|A| + |B|$.
	First, suppose $|A| + |B| = 4$, which is the base case.
	By assumption, there exists $u \in \Sigma^*$
	such that $3 \le |\applytrans{\delta}{A}{u}| + |\applytrans{\delta}{B}{u}|  < 4$, which
	implies that one of both sets must be a singleton and the other a $2$-set.
	Hence, $u \in \Sigma^*$ distinguishes both subsets.
	Next, suppose $|A| + |B| > 4$. Assume all pairs of distinct subsets $A', B' \subseteq Q$
	with $4 \le |A'| + |B'| < |A| + |B|$
	and $|A'|, |B'| \ge 2$ are distinguishable.
	By the assumption from~\cref{dist:item:dcps} there exists
	$u \in \Sigma^*$
	such that  $3 \le |\applytrans{\delta}{A}{u}| + |\applytrans{\delta}{B}{u}|  < |A| + |B|$
	and $\applytrans{\delta}{A}{u} \ne \applytrans{\delta}{B}{u}$.
	If one of both subsets is a singleton, the other subset has size at least two.
	So, $u$ distinguishes~$A$ and~$B$.
	Otherwise, as $\applytrans{\delta}{A}{u} \ne \applytrans{\delta}{B}{u}$, by induction hypothesis, 
	there exists $v \in \Sigma^*$
	distinguishing $\applytrans{\delta}{A}{u}$ from $\applytrans{\delta}{B}{u}$.
	Then~$uv$ distinguishes~$A$ from~$B$.
\end{proof}

\Crefitem{sync_max:thm:char_distinguishable}{dist:item:dcps} is useful for an inductive reasoning
because of the fact that it only demands to compress the subsets (while retaining their distinctiveness)
instead of instantly mapping precisely one to a singleton subset.

\begin{remark}
 By \cref{sync_max:ex:distinguishable-trap} 	
 distinguishability of all pairs of distinct $(n-1)$-sets 
 is not sufficient for complete distinguishability.
 With \cref{sync_max:prop:ksets},
 this also implies that there exists no~$k$
 such that distinguishability of all pairs of distinct $k$-sets
 implies complete distinguishability.
\end{remark}

\begin{remark}
  Having for each $q \in Q$ a word $u_q \in \Sigma^*$
  that distinguishes $Q$ form $Q \setminus \{q\}$
  is necessary, and, by \crefitem{sync_max:thm:char_distinguishable}{dist:item:single_point}
  also sufficient, for complete distinguishability.
  Let $p,q,r \in Q$ be distinct.
  Assume $\delta(r, u_q) = \delta(r, u_p)$.
  Then we have $\delta(p, u_q) = \delta(r, u_q) = \delta(r, u_p) = \delta(q, u_p)$.
  Since $\delta(q, u_q) \ne \delta(p, u_q)$,
  we find  $\delta(q, u_q) \ne \delta(q, u_p)$.
  Thus we conclude that each word $u_q$ induces a different
  transformation on the states.
  Hence, the transformation monoid of completely distinguishable
  automata with at least three states contains at least~$|Q|$ different elements.
\end{remark}

\begin{corollary}\label{sync_max:cor:sufficient-compl-dist}
	Let $\semiautomat{A}{Q}{\Sigma}{\delta}$.
	If for all $r \in Q$
	and $p, q \in Q \setminus \{q\}$
	there exists $u \in \Sigma^*$
	such that $\delta(p, u) = \delta(q, u)$, $\delta(r, u) = r$,
	and $\delta(s, u) \ne r$ for all $s \in Q \setminus \{r\}$, then $\mathcal A$
	is completely distinguishable.
\end{corollary}
\begin{proof}
	This follows directly from \crefitem{sync_max:thm:char_distinguishable}{dist:item:pqr-r-one-lane}
	since it equals this statement in the special case that the word fixes the state~$r$.
\end{proof}

\begin{corollary}\label{cor:compl-dist-transitive}
	Let $\semiautomat{A}{Q}{\Sigma}{\delta}$
	be an automaton such that the group of units of~$\TM{\mathcal A}$
	is transitive.
	Then $\mathcal A$ is completely distinguishable
	if and only if $Q$ is distinguishable from $Q \setminus \{q\}$
	for some $q \in Q$.		
\end{corollary}
\begin{proof}
	First, suppose $Q$ and $Q \setminus \{q\}$ are distinguishable for $q \in Q$.
	Then $|\delta(Q \setminus \{q\}, u)| = 1$
	with $\delta(q, u) \notin \delta(Q \setminus\{q\}, u)$
	for some $u \in \Sigma^*$.
	Let $p \in Q$.
	Since~$\TM{\mathcal A}$ contains a transitive permutation group, 
	there exists $v \in \Sigma^*$ inducing a permutation on $Q$
	such that $\delta(p, v) = q$.
	Then we have $\applytrans{\delta}{Q \setminus \{p\}}{vu} = \applytrans{\delta}{Q\setminus\{q\}}{u}$
	and $\applytrans{\delta}{p}{vu} = \applytrans{\delta}{q}{u}$.
	So $\applytrans{\delta}{p}{vu} \notin \applytrans{\delta}{Q \setminus\{p\}}{vu}$
	and $|\applytrans{\delta}{Q\setminus \{p\}}{vu}| = 1$.
	By~\crefitem{sync_max:thm:char_distinguishable}{dist:item:single_point},
	it follows that $\mathcal A$ is completely distinguishable.	
	The converse implication is clear.
\end{proof}

\begin{example}  \label{sync_max:ex:compl-dist-non-cr}
	Consider the automaton
	in \cref{sync_max:fig:compl_dist_not_compl_reach:1}.
	The letter $a$ distinguishes $\{0,1,2,3\}$
	from the whole state set. 
	Since the letter $b$ generates a transitive permutation group,
	the automaton is completely distinguishable by \cref{cor:compl-dist-transitive}.	
	The automaton is not completely reachable
	since it does not have a letter of defect one.
\end{example}

\begin{example} \label{sync_max:ex:compl-dist-non-cr2}
	\Cref{sync_max:fig:compl_dist_not_compl_reach:2}
	depicts an automaton that is easily seen
	to be completely distinguishable
	and that is not completely reachable
	since it is not strongly connected.
\end{example}

\begin{example} \label{syncmax:ex:compl-dist-not-SCC}
	The word $ab$ distinguishes $\{1,2,3\}$ and $\{4\}$.
	The word $ecab$ distinguishes $\{2,3,4\}$ and $\{1\}$.
	The word $bc$ distinguishes $\{1,2,4\}$ and $\{3\}$.
	The word $ac$ distinguishes $\{1,3,4\}$ and $\{2\}$.
	So, by \crefitem{sync_max:thm:char_distinguishable}{dist:item:single_point},
	the automaton from \cref{sync_max:fig:compl_dist_example} is completely distinguishable,
	but it is not completely reachable since it is not strongly connected.
	Note that by adding a simple idempotent mapping $1$ to $2$, we obtain
	a strongly connected SI-automaton that is completely distinguishable,
	but whose automaton graph is not strongly $\frac{3}{2}$-connected
	since all paths from $2$ to $4$ have to pass through~$1$.
\end{example}

\begin{figure}[ht]
	\centering 
	\begin{subfigure}[t]{0.37\textwidth}
		\centering
		\begin{tikzpicture}[shorten >=1pt,->,node distance=4cm]
				\foreach \name/\angle/\text in {P-n-2/234/3, P-n-1/162/4, 
					P-0/90/0, P-1/18/1, P-2/-54/2}
				\node[state] (\name) at (\angle:7em) {$\text$};
				
				\path[->] (P-n-2) edge [left] node {$b$} (P-n-1);
				\path[->] (P-n-1) edge [above,pos=0.35] node {$b$} (P-0);
				\path[->] (P-0) edge [above] node {$b$} (P-1);
				\path[->] (P-1) edge [right] node {$b$} (P-2);
				\path[->] (P-2) edge [below] node {$a,b$} (P-n-2);
				\path[->] (P-0)   edge [right] node {$a$} (P-n-2)
		                  (P-1)   edge [above] node {$a$} (P-n-2);					
			\end{tikzpicture}
		\caption{The Strongly connected automaton from Example~\ref{sync_max:ex:compl-dist-non-cr}.}
		\label{sync_max:fig:compl_dist_not_compl_reach:1}
	\end{subfigure}\hfill
	\begin{subfigure}[t]{0.29\textwidth}
		\centering		
		\begin{tikzpicture}
			\node[state] (1) {$1$};
			\node[state, above = of 1] (2) {$2$};
			\node[state, left  = of 2] (3) {$3$};    		
			
			\path[->] (3) edge [left]  node {$a$} (1)
			          (2) edge [right] node {$b$} (1)
			          (2) edge [above] node {$c$} (3);
		\end{tikzpicture}		
		\caption{The automaton from Example~\ref{sync_max:ex:compl-dist-non-cr2}.} 
		\label{sync_max:fig:compl_dist_not_compl_reach:2}
	\end{subfigure}\hfill
    \begin{subfigure}[t]{0.29\textwidth}
    	\centering
    		\begin{tikzpicture}[node distance = 3.5em and 1.5em]
    		\node[state] (1) {$1$};
    		\node[state, above right = of 1] (2) {$2$};
    		\node[state, above left  = of 1] (3) {$3$};
    		\node[state, below = of 1]       (4) {$4$};				
    		
    		\path[->] (3) edge [left]  node {$a$} (1)
    		          (2) edge [right] node {$b$} (1)
    		          (1) edge [bend left=10, right] node {$c$} (4)
    		          (4) edge [bend left=10, left]  node {$d$} (1)
    		          (4) edge [bend left, left]  node {$e$} (3);
    	\end{tikzpicture}
        \caption{The automaton from \cref{syncmax:ex:compl-dist-not-SCC}}     	
        \label{sync_max:fig:compl_dist_example}    
    \end{subfigure}	\caption[Completely distinguishable, but not completely reachable]{Completely 
    distinguishable automata that are not completely reachable. Self-loops omitted.}
	\label{sync_max:fig:compl_dist_not_compl_reach}
\end{figure}

It is easily seen that quotient automata of completely reachable automata
are completely reachable. However, complete distinguishability is not preserved
under taking quotient automata as shown in \cref{sync_max:fig:compl_dist_not_compl_reach} next.

% Und sync-maximal erhalten unter Quotientautomaten?
\begin{example} \label{sync_max:fig:compl_dist_not_compl_reach}
	Automata  $\semiautomat{A}{Q}{\Sigma}{\delta}$ as in \cref{sync_max:ex:compl-dist-non-cr}
	also show that complete distinguishability
	is not preserved when taking quotient automata.
	For consider the obvious generalization	where the letter~$b$ cyclically permutes $n$ states
	and suppose we have some non-trivial $\TS{\mathcal A}$-congruence $\sim$.
	%Let $p \in Q$ be the state fixed by~$a$.
	Let $s \in Q$ be the state fixed by~$a$.
	Since the letter $b$ cyclically permutes the whole state set,
	all congruence classes for~$\sim$ have the same size
	and we find $t \in Q$ with $s \ne t$
	such that $s \sim t$.
		
	First, we show that for every $u \in \Sigma^*$ not permuting the set of states
	we have $\delta(p, u) \sim \delta(q, u)$
	for all $p,q \in Q$, i.e., the letter~$a$ induces a constant in $\mathcal A / {\sim}$).
	Let $p,q \in Q$. We can assume $p \not\sim q$.
	For $u \in \Sigma^*$ not permuting~$Q$
	we can write $u = b^iav$ for some $v \in \Sigma^*$ and $i \ge 0$.
	If $\delta(p, b^i) \ne s$ 
	and $\delta(q, b^i) \ne s$,
	then $\delta(p, b^ia) = \delta(q, b^ia)$
	and so $\delta(p, u) = \delta(q, u)$.
	Otherwise, without loss of generality, assume $\delta(p, b^i) = s$.
	Then $\delta(q, b^i) \ne s$.
	Now, we have $\delta(p, b^ia) = s$
	and $\delta(q, b^ia) = \delta(t, a)$.
	Furthermore, $s \sim t$
	implies $\delta(p, b^ia) = s = \delta(s, a) \sim \delta(t, a) = \delta(q, b^ia)$
	and thus $\delta(p, u) \sim \delta(q, u)$.
		
	Next, assume we want to distinguish $Q / {\sim}$
	and $Q / {\sim} \setminus \{ [p]_{\sim} \}$
	in $\mathcal A / {\sim}$ for some $p \in Q$.
	Then a distinguishing word $u$
	maps all $q \in Q$ with $q \not\sim p$
	into an equivalence class
	and for all $r \in Q$ with $r \sim p$
	we have $\delta(r, u) \not\sim \delta(q, u)$ for $q \in Q$ with $q \not\sim p$.
	However, since a distinguishing word cannot permute
	the whole state set, as shown above this is not possible.	
\end{example}

Next, we show that a binary completely distinguishable automaton with a letter of rank two, such as the one from \cref{sync_max:fig:compl_dist_not_compl_reach:1}, must contain a letter cyclically permuting the states.

\begin{proposition} \label{syncmax:prop:binary-rank-two-letter}
	Let $\semiautomat{A}{Q}{\{a,b\}}{\delta}$ be a binary automaton with at least three states
	where~$a$ has rank two.
	If $\mathcal A$ is completely distinguishable,
	then~$b$ cyclically permutes the states.
\end{proposition}
\begin{proof}
	Assume $\delta(p_1, b) = \delta(p_2, b)$
	for two distinct states $p_1, p_2 \in Q$.
	Then a word distinguishing $\{p_1\}$ from $Q \setminus\{p_1\}$
	must start with~$a$.
	This implies $\delta(p_1, a) \notin \delta(Q \setminus \{p_1\}, a)$.
	Since~$a$ has rank two,
	we get $|\delta(Q \setminus \{p_1\}, a)| = 1$
	and so~$a$ distinguishes both subsets.
	By a similar reasoning with $\{p_2\}$ and $Q \setminus \{p_2\}$,
	we find $|\delta(Q \setminus \{p_2\}, a)| = 1$.
	Since~$|Q| \ge 3$,
	this yields that $Q \setminus \{p_1\}$
	and $Q \setminus \{p_2\}$ are mapped to the same
	state and so $\delta(p_1, a) = \delta(p_2, a)$,
	giving a contradiction
	to $\delta(p_1, a) \notin \delta(Q \setminus \{p_1\}, a)$.
	Thus, no two distinct states are mapped
	to the same state by~$b$,
	and by finiteness we find that~$b$
	permutes the states.
	
	Next, by assumption we have $Q = A \cup B$
	where $A = \delta^{-1}(p, a)$
	and $B = \delta(q, a)$ for some $p,q \in Q$.
	Assume there exist two distinct $|A|, |B| \ge 2$.
	Each word distinguishing two subsets
	must have the form $b^ia$ for some $i \ge 0$,
	since~$a$ has rank two.
	We have $\delta^{-1}(p, b^ia) = \delta^{-1}(A, b^i)$
	and $\delta^{-1}(q, b^ia) = \delta^{-1}(B, b^i)$
	and the two subsets have size at least two
	and partition~$Q$ since~$b$ permutes the states.
	However, a word distinguishes a state~$s$
	from $Q \setminus \{s\}$
	if and only if $\{s\}$
	and $Q \setminus \{s\}$
	are the $\delta^{-1}$-images of two states.
	Thus, one of both subsets $A$ or $B$
	must be a singleton.
	Without loss of generality, let~$A = \{p_1\}$ and consider some state $p_2 \in Q$.
	A word $b^ia$ distinguishes $p_2$ and $Q \setminus \{p_2\}$
	if and only if $p_2$ is mapped to~$p_1$
	by~$b^i$. Thus we find that~$b$ cyclically permutes the states.	
\end{proof}

Lastly, we also note a readily seen equivalence when the whole state set is distinguishable
from a set of the form $Q \setminus\{q\}$, since these sets appear in \crefitem{sync_max:thm:char_distinguishable}{dist:item:single_point}

% in beweis oben nutzen?
\begin{proposition} \label{sync_max:prop:dist-n-1-set}
	Let $\semiautomat{A}{Q}{\Sigma}{\delta}$. Then the following are equivalent.
	\begin{enumerate}
		\item $Q \setminus \{q\}$ is distinguishable from~$Q$.
		
		\item There exist $u_1, \ldots, u_m \in \Sigma$
		such that for  $A_i = \delta(Q \setminus \{q\}, u_1 \cdots u_i)$
		and $q_i = \delta(q, u_1 \cdots u_i)$ where $i \in \{1,\ldots,m\}$
		we have  $|A_m| = 1$
		and $q_i \notin A_i$ for all $i \in \{1,\ldots,m\}$.
	\end{enumerate}
\end{proposition}

\begin{example}
	Let $\semiautomat{A}{\{1,2,3,4\}}{\Sigma}{\delta}$
	with 
	\[
	\Sigma = \{ a \mid \exists y_1, y_2 \in Q \colon y_1 \ne y_2 \land
	|\delta^{-1}(y_1, a)| = |\delta^{-1}(y_1, a)| = 2 \}
	\]
	Then disjoint pairs of $2$-sets are indistinguishable
	and non-disjoint pairs of $2$-sets are distinguishable.
	Furthermore, two sets of size at least three are indistinguishable.
	A set of size two is distinguishable from any set of size at least three.	
\end{example}

\subsection{Further Examples}
\label{sync_max:sec:furtherexamples} 

%!TEX root = ../../springer-mono-format/book.tex
%!TEX root = ../../thesis-trier.tex

\Cref{sync_max:ex:2-sets,sync_max:ex:2-sets-binary} gave infinite families of automata that are not completely distinguishable (as they are not synchronizing and have more than three states), but in which all $2$-sets are pairwise distinguishable.
Furthermore, in the synchronizing binary SI-automaton with trap state $t$ from \cref{prelim:fig:binarySI-sync-trap}, every pair of distinct $2$-sets is distinguishable, whereas any two subsets of the form $A$ and $A \cup \{t\}$, where $t \notin A$ and $|A| \ge 2$, are indistinguishable.
Here, in \cref{sync_max:ex:2-sets-not-all-dist,sync_max:ex:2-sets-but-not-all-new}, we give further automata that are not completely distinguishable, but are synchronizing and have the property that all distinct pairs of $2$-sets are distinguishable.
The automaton from \cref{sync_max:ex:2-sets-not-all-dist} 
is even a completely reachable SI-automaton.

\begin{example}\label{sync_max:ex:2-sets-not-all-dist}
	\Cref{sync_max:fig:SCC_not_sufficient} gives a synchronizing and strongly connected SI-automaton
	with the following properties.
	\begin{enumerate}		
		\item A set $\{q_1, \ldots, q_m\}$ with $q_1 < \cdots < q_m$
		is indistinguishable from $\{q_1, q_m\}$.
				
		\item All distinct pairs of $2$-subsets of states are distinguishable.
		
		\item For $A, B \subseteq Q$ with $|A|,|B| \ge 2$ there exists $u \in \Sigma^*$
		such that $|\applytrans{\delta}{A}{u}| + |\applytrans{\delta}{B}{u}| < |A| + |B|$.
		
		\item For $A, B \subseteq Q$ with $|A|,|B| \ge 2$ and $A \ne B$ there exists $u \in \Sigma^*$
		such that $|\applytrans{\delta}{A}{u}| \ne |\applytrans{\delta}{B}{u}|$.
	\end{enumerate}
	The automaton shows that the latter two 
	properties, which are similar to \crefitem{sync_max:thm:char_distinguishable}{dist:item:dcps},
	are not sufficient for complete distinguishability.	
	It is interesting to note that adding a simple idempotent letter mapping $1$ to $n$ to the automaton in \cref{sync_max:fig:sync-two-sets-dist} yields a completely distinguishable automaton,
	which follows from \cref{sync_max:thm:IS_32SCC} since the resulting automaton graph 
	is strongly $\frac{3}{2}$-connected.
	We note that the automaton from \cref{sync_max:ex:2-sets} 
	also has the property that two distinct subsets of size at least two can be mapped
	to subsets of different sizes, but it is not even synchronizing.
		
	\begin{figure}[htb]
		\centering
		\begin{tikzpicture}[inner sep=1pt,outer sep=2pt] 
			\node[state] (1) {$1$};
			\node[state, right = of 1] (2) {$2$};
			\node[state, right = of 2] (3) {$3$};
			\node[state, right = of 3] (4) {$4$};	
			\node[right = of 4] (dots) {$\cdots$};
			\node[state, right = of dots] (n1) {$n-1$};
			\node[state, right = of n1] (n) {$n$};
			
			\path[->] (1) edge [bend left, above] node {$a_1$} (2)
			(2) edge [bend left, above] node {$a_2$} (3)
			(3) edge [bend left, above] node {$a_3$} (4)
			(4) edge [bend left, above] node {$a_4$} (dots);
			
			\path[->] (dots) edge [bend left, above] node {$a_{n-2}$} (n1)
			(n1)   edge [bend left, above] node {$a_{n-1}$} (n);
			
			\path[->] (n) edge [bend left, below] node {$b_{n}$} (n1)
			(n1) edge [bend left, below] node {$b_{n-1}$} (dots)
			(dots) edge [bend left, below] node {$b_{n-2}$} (4)
			(4) edge [bend left, below] node {$b_{4}$} (3)
			(3) edge [bend left, below] node {$b_3$} (2)
			(2) edge [bend left, below] node {$b_2$} (1);
		\end{tikzpicture}
		\caption[Pairs of $2$-sets distinguishable but not completely distinguishable]{Strongly connected and synchronizing SI-automata (self-loops omitted) where 
			all pairs of $2$-sets are distinguishable
			and a set $\{q_1, \ldots, q_m\}$ with $q_1 < \cdots < q_m$
			is indistinguishable from $\{q_1, q_m\}$. See \cref{sync_max:ex:2-sets-not-all-dist}.}	
		\label{sync_max:fig:sync-two-sets-dist}
		\label{sync_max:fig:SCC_not_sufficient}
	\end{figure}
\end{example}

\begin{example}\label{sync_max:ex:2-sets-but-not-all-new}
	Let $\semiautomat{A}{\{1,2,3,4\}}{\{a,b,c,d,e\}}{\delta}$ 
	be the automaton from \cref{sync_max:fig:2-sets-but-not-all:abcde}
	over the alphabet $\{a,b,c,d,e\}$.
	\Cref{sync_max:fig:2-sets-but-not-all:abcd} and \cref{sync_max:fig:2-sets-but-not-all:left} 
	show $\mathcal A_{|\{a,b,c,d\}}$
	and $\mathcal A_{|\{a,b,c\}}$, respectively.
	
	\begin{figure}[htb]
		\centering 
		\begin{subfigure}{0.33\textwidth}
			\begin{center} 
				\begin{tikzpicture}[node distance = 1.45cm and 1cm]
					\node[state] (1) {$1$};
					\node[state, right = of 1] (2) {$2$};
					\node[state, below = of 2] (3) {$3$};
					\node[state, below = of 1] (4) {$4$};
					
					\path[->] (1) edge [loop above] node {$a,c, d$} (1)
					(3) edge [right, align=left]  node {$a, b$ \\ $c, d$} (2)
					(4) edge [left, align=right] node {$a, b$ \\ $c, e$} (1)
					(4) edge [below] node {$d$} (2);
					
					\path[->] (1) edge [bend left,above] node {$b, e$} (2)
					(2) edge [bend left,above] node {$a,b$} (1)
					(2) edge [loop above] node {$c,d, e$} (2);			
					
					\path[->] (3) edge [below] node {$e$} (4);	
				\end{tikzpicture}
			\end{center}
			\caption{$\{1,2,3,4\}$ and $\{1,3,4\}$ are indistinguishable.}
			\label{sync_max:fig:2-sets-but-not-all:abcde}
		\end{subfigure}
		\hfill 
		\begin{subfigure}{0.3\textwidth}
			\begin{center} 
				\begin{tikzpicture}[node distance = 1.45cm and 1cm]
					\node[state] (1) {$1$};
					\node[state, right = of 1] (2) {$2$};
					\node[state, below = of 2] (3) {$3$};
					\node[state, below = of 1] (4) {$4$};
					
					\path[->] (1) edge [loop above] node {$a,c, d$} (1)
					(3) edge [right, align=left]  node {$a, b$ \\ $c, d$} (2)
					(4) edge [left, align=right] node {$a, b$ \\ $c$} (1)
					(4) edge [below] node {$d$} (2);
					
					\path[->] (1) edge [bend left,above] node {$b$} (2)
					(2) edge [bend left,above] node {$a,b$} (1)
					(2) edge [loop above] node {$c,d$} (2);
				\end{tikzpicture}
			\end{center}
			\caption{$\{1,2,3,4\}$ and $\{1,2,3\}$ are indistinguishable.}
			\label{sync_max:fig:2-sets-but-not-all:abcd}
		\end{subfigure}%
		\hfill	
		\begin{subfigure}{0.3\textwidth}
			\begin{center} 
				\begin{tikzpicture}[node distance = 1.45cm and 1cm]
					\node[state] (1) {$1$};
					\node[state, right = of 1] (2) {$2$};
					\node[state, below  = of 2] (3) {$3$};
					\node[state, below  = of 1] (4) {$4$};
					
					\path[->] (1) edge [loop above] node {$a,c$} (1)
					(3) edge [right, align=left]  node {$a, b$ \\ $c$} (2)
					(4) edge [left,  align=right] node {$a,b$ \\ $c$} (1);
					
					\path[->] (1) edge [bend left,above] node {$b$} (2)
					(2) edge [bend left,above] node {$a,b$} (1)
					(2) edge [loop above] node {$c$} (2);
				\end{tikzpicture}
			\end{center}
			\caption{$\{1,2,3\}$ and $\{3,4\}$ are indistinguishable.}
			\label{sync_max:fig:2-sets-but-not-all:left}
		\end{subfigure}%
		\caption[Non-completely distinguishable but pairs of $2$-sets distinguishable]{Automata that are not completely distinguishable, but in which all pairs of distinct $2$-sets are distinguishable; see \cref{sync_max:ex:2-sets-but-not-all-new}.}
		\label{sync_max:fig:2-sets-but-not-all}
	\end{figure}
	
	% https://www.reddit.com/r/LaTeX/comments/jwdp5y/table_too_wide_and_too_long_to_fit_in_one_page/
	\begin{table}
		\caption[Distinguishing words for pairs of subsets]{Pairs of subsets
			in the automaton from \cref{sync_max:fig:2-sets-but-not-all:abcde}
			that are distinguishable by the indicated word,
			or indistinguishable as indicated by a hyphen.}
		\resizebox{\textwidth}{!}{
			$
			\begin{array}{r|cccccccccccccccccc}
				& \{1,2\} & \{1,3\} & \{1,4\} & \{2,3\} & \{2,4\} & \{3,4\} & \{1,2,3\} & \{1,2,4\} & \{1,3,4\} & \{2,3,4\}  \\ \hline
				\{1,3\}     & a       & -       &         &         &         &         &           &           &           &            \\ 
				\{1,4\}     & c       & a       & -       &         &         &         &           &           &           &            \\
				\{2,3\}     & a       & b       & a       & -       &         &         &           &           &           &             \\
				\{2,4\}     & b       & a       & c       & a       & -       &         &           &           &           &             \\
				\{3,4\}     & a       & b       & a       & c       & a       & -       &  \\
				\{1,2,3\}   & a       & b       & a       & c       & b       & d       & -         &           &           &   \\
				\{1,2,4\}   & e       & b       & c       & c       & b       & d       & a         & -         &           &   \\ 
				\{1,3,4\}   & a       & b       & a       & c       & b       & d       & ee        & a         & - \\
				\{2,3,4\}   & a       & b       & a       & c       & b       & ec      & d         & a, d      & d         & -       \\
				\{1,2,3,4\} & a       & b       & a       & c       & b       & d       & eee       & a         & -         & d
			\end{array} 
			$
		}	
		\label{syncmax:tab:2-set-but-not-all}
	\end{table}	
	
	These automata are not completely distinguishable.\footnote{However, similarly
		as in \cref{sync_max:ex:2-sets-not-all-dist}, for each pair of distinct $A,B$ with $|A|,|B| \ge 2$
		there exists $u \in \Sigma^*$ such that $|\delta(A, u)| \ne |\delta(B, u)|$.}
	\Cref{syncmax:tab:2-set-but-not-all}
	shows distinguishing words for all distinguishable pairs in $\mathcal A$.	
	We have the
	following maximal classes of pairwise indistinguishable
	subsets of states:
	\begin{align*}
		\mathcal A                & : \{\{1,3,4\}, \{1,2,3,4\}\}, \\ 
		\mathcal A_{|\{a,b,c,d\}} & : \{\{1,2\},\{1,2,4\}\}, \{ \{3,4\}, \{2,3,4\}\}, \{ \{1,2,3\}, \{1,3,4\}, \{1,2,3,4\}\}, \\
		\mathcal A_{|\{a,b,c\}}   & : \{\{1,2\},\{1,2,4\}\}, \{ \{3,4\},  \{2,3,4\}, \{1,2,3\}, \{1,3,4\}, \{1,2,3,4\}\}.  	
	\end{align*}
 	For $\mathcal A_{|\{a,b,c\}}$ and $\mathcal A_{\{a,b,c,d\}}$ this follows from the following arguments:
	\begin{enumerate} 
	\item \label{sync_max:ex:2-sets-but-not-all-new:1} We have  $\delta^{-1}(1, a) = \{1,2,4\}$.
	So, for two distinct $2$-sets $A, B$ of states
	such that $A \subseteq \{1,2,4\}$ and $3 \in B$
	the letter $a$ yields a distinguishing word.
	As $\delta^{-1}(1, c) = \{1,4\}$
	and $\delta^{-1}(1, b) = \{2,4\}$
	every pair of distinct $2$-sets from $\{1,2,4\}$
	is distinguishable
	and as $\delta^{-1}(2,c) = \{2,3\}$
	and $\delta^{-1}(2, b) = \{1,3\}$
	every pair of distinct $2$-sets $A, B$ with $3 \in A \cap B$
	is distinguishable.
	So, we find that all pairs of distinct $2$-sets
	are distinguishable in $\mathcal A_{\{a,b,c\}}$.
	
	We also obtain that every $2$-set not equal to $\{3,4\}$
	is distinguishable from $\{1,2,3,4\}$, and that $\{1,2,4\}$
	is distinguishable from every other $3$-set.
		
	\item \label{sync_max:ex:2-sets-but-not-all-new:2} Let $x \in \{a,b,c\}$. For each $A$ with $|A| \ge 3$ and $A \ne \{1,2,4\}$ 
	we have 
	\[ 
	\delta(A, x) = \delta(\{3,4\}, x) = \{1,2\}.
	\]
	Therefore, any two sets $A,B$ with $|A|,|B|\ge 3$ and $A,B \ne \{1,2,3\}$
	are indistinguishable in $\mathcal A_{|\{a,b,c\}}$.
	Moreover, $\{3,4\}$ is indistinguishable
	from every subsets $A$ with $|A| \ge 3$ and $A \ne \{1,2,4\}$.
	Since $\delta(\{1,2\}, x) = \delta(\{1,2,4\}, x)$,
	the sets $\{1,2\}$ and $\{1,2,4\}$ are also indistinguishable in $\mathcal A_{|\{a,b,c\}}$.
	
	\Cref{sync_max:ex:2-sets-but-not-all-new:1,sync_max:ex:2-sets-but-not-all-new:2} yield that
	\[
	\{\{1,2\},\{1,2,4\}\} \text{ and } \{ \{3,4\},  \{2,3,4\}, \{1,2,3\}, \{1,3,4\}, \{1,2,3,4\}\}
	\]
	are two maximal classes of pairwise indistinguishable subsets in $\mathcal A_{|\{a,b,c\}}$.	
		
	\item We have $\delta^{-1}(2, d) = \{2,3,4\}$.
	Hence, every subset of $\{2,3,4\}$
	is distinguishable in $\mathcal A_{|\{a,b,c,d\}}$
	from every subset $A$ with $|A| \ge 2$
	and $1 \in A$. Since $\delta(\{3,4\}, d) = \delta(\{2,3,4\}, d) = \{2\}$,
	\[
	\delta(\{1,2,3\}, d) = \delta(\{1,2,4\}, d) = \delta(\{1,3,4\}, d) = \delta(\{1,2,3,4\}, d) = \{1,2\},
	\]
	and $\{ \{3,4\},  \{2,3,4\}, \{1,2,3\}, \{1,3,4\}, \{1,2,3,4\}\}$
	is a maximal set of pairwise indistinguishable subsets in $\mathcal A_{|\{a,b,c\}}$,
	we find that 
	\[
	\{ \{3,4\}, \{2,3,4\}\} \text{ and }
	\{ \{1,2,3\}, \{1,3,4\}, \{1,2,3,4\}\}
	\]
	are maximal sets of pairwise indistinguishable subsets in $\mathcal A_{|\{a,b,c,d\}}$.
	\end{enumerate}
\end{example}

Recall that distinguishability of pairs of distinct $2$-sets
does not imply synchronizability, see \cref{sync_max:fig:non-sync-two-sets-dist}.
However, this implication holds for strongly connected
automata by \cref{sync_max:prop:2-sets-dist-SCC-implies-sync}.
